% Generated draft: requires mathematical and source review.
\documentclass[11pt]{article}
\usepackage[margin=25mm]{geometry}
\usepackage{fontspec}
\setmainfont{Latin Modern Roman}
\usepackage{amsmath,amssymb,mathtools,unicode-math}
\setmathfont{Latin Modern Math}
\usepackage{enumitem,needspace,xurl,hyperref}
\setlength{\parindent}{0pt}
\setlength{\parskip}{5pt}
\setlength{\emergencystretch}{3em}
\newcommand{\N}{\mathbb{N}}
\newcommand{\Z}{\mathbb{Z}}
\newcommand{\Q}{\mathbb{Q}}
\newcommand{\R}{\mathbb{R}}
\newcommand{\C}{\mathbb{C}}
\newcommand{\F}{\mathbb{F}}
\newcommand{\A}{\mathbb{A}}
\newcommand{\PP}{\mathbb{P}}
\newcommand{\E}{\mathbb{E}}
\newcommand{\eps}{\varepsilon}
\newcommand{\dd}{\,\mathrm{d}}
\begin{document}
\section*{UnsolvedMath ID: 800713}
\textbf{Source problem code:} catalogue entry
\subsection{Definitions and mathematical statement}
Let \(K\) be a global field of characteristic \(p\), as stipulated in the source, and let \(A\) be an abelian variety over \(K\). Write \(\hat A\) for its dual abelian variety and set
\[
r=\operatorname{rank}_{\Z} A(K),\qquad
M=A(K)/A(K)_{\mathrm{tors}},\qquad
\hat M=\hat A(K)/\hat A(K)_{\mathrm{tors}}.
\]
Here the subscript \(\mathrm{tors}\) denotes the torsion subgroup. For each place \(v\) of \(K\), let \(K_v\) be the completion. Interpreting the source's notation \(\mathrm{III}(A)\) in its usual BSD sense, define the Tate--Shafarevich group by
\[
\mathrm{III}(A/K)=\ker\!\left(H^1(K,A)\longrightarrow\prod_v H^1(K_v,A)\right),
\]
where \(H^1\) denotes Galois cohomology.

The canonical N\'eron--Tate pairing has the form
\[
h_A:M\times\hat M\longrightarrow\R.
\]
For integral bases \(P_1,\ldots,P_r\) of \(M\) and \(Q_1,\ldots,Q_r\) of \(\hat M\), its regulator is
\[
R_A=\left|\det\bigl(h_A(P_i,Q_j)\bigr)_{1\leq i,j\leq r}\right|,
\]
with the empty determinant equal to \(1\) when \(r=0\). This is the conventional interpretation of the determinant in the displayed source formula. The source instead describes \(h_A\) simply as a pairing on \(A(K)\) modulo torsion; it supplies no polarization or identification with the dual, so this interpretation requires confirmation.

Let \(L^*(A,s)\) denote the source's \(L\)-function with respect to a ``suitably normalized Tamagawa measure.'' The supplied BSD conjecture asserts analytic continuation of \(L^*(A,s)\) to the complex plane and the leading-coefficient identity
\[
\lim_{s\to1}\frac{L^*(A,s)}{(s-1)^r}
=
\frac{|\mathrm{III}(A/K)|\,R_A}
{|A(K)_{\mathrm{tors}}|\,|\hat A(K)_{\mathrm{tors}}|}.
\]
For this expression to be a finite numerical identity, finiteness of \(\mathrm{III}(A/K)\) is required; the supplied statement does not separately assert that finiteness. The source does not specify the normalization of \(L^*\), the Tamagawa measure, or the height pairing sufficiently to determine the numerical identity independently. It also does not specify whether continuation to the complex plane means holomorphic or meromorphic continuation. These qualifications remain unresolved.
\subsection{Short English statement}
For every abelian variety over a global field of the stated characteristic, the supplied BSD conjecture predicts continuation of its normalized \(L\)-function and a leading coefficient at \(s=1\) determined by the Tate--Shafarevich group, a height regulator, and rational torsion groups. The normalization, pairing convention, continuation type, and required finiteness need clarification.
\subsection{Sources}
\begin{itemize}
\item Local supplied UnsolvedMath registry record, permanent ID \(800713\), exact source problem code \texttt{catalogue entry}; upstream ID \(800713\), revision \texttt{d371ee917ac9647fbb06cd6fc529deec6f3dda82}. The record was created and updated on \texttt{2026-10-16T14:00:00Z} and labels the problem \emph{open}; this is supplied registry status, not an independently verified literature assessment.
\item Auel, Brussel, Garibaldi, Vishne, \emph{Open Problems on Central Simple Algebras} (2010), Conjecture 27, identified in the record as the BSD Conjecture. Supplied source URL: \url{https://arxiv.org/abs/1006.3304}. The registry attributes the item to these authors and gives proposed year 2010; its background identifies the context as the BSD Conjecture.
\item The supplied provenance records access on \texttt{2026-10-16} and extraction method \texttt{arxiv\_latex}. These are supplied provenance statements; the source was not independently visited or verified for this draft.
\end{itemize}
\end{document}
